% -----------------------------------------------------------------------------
% Feynman-diagram combinatorics for beginners, strict SI conventions
% -----------------------------------------------------------------------------

Wick 定理把时序乘积变成所有可能的缩并，而图的数值系数还取决于相同缩并的计数、费米算符重排的奇偶性和图的自同构。几个最小例子的算符重排可以把这些组合因子的来源逐项看清。场的物理作用量仍采用 SI 单位；截去外腿以后，图上使用
\eqnrefs{eq:si-reduced-propagators-r10,eq:si-reduced-vertices-r10}
定义的约化因子。

给定已经由 Wick 定理生成的图，沿费米子-数 arrow 连续读取：入射电子给 $u$，出射电子给 $\bar u$；入射正电子给 $\bar v$，出射正电子给 $v$；每经过一个相互作用顶角写相应 Dirac 矩阵，每经过一条内部费米子线写
\eqnrefs{eq:si-reduced-propagators-r10} 的约化传播子。若有多条不等价的费米收缩，先按固定外部算符次序分别写振幅，再由置换宇称决定相对号。


Feynman 规则可还原为算符代数：顶角系数来自相互作用作用量与等价收缩的组合数；内部线来自收缩；对称性因子来自保持外部标签不变的图自同构；全同外部费米子与闭合费米子圈的负号都来自反对易关系。图规则只是这些算符记账步骤的压缩表示。

